\documentclass{amsart}
% For dealing with references we use the comment environment
\usepackage{verbatim}
\newenvironment{reference}{\comment}{\endcomment}
%\newenvironment{reference}{}{}
\newenvironment{slogan}{\comment}{\endcomment}
\newenvironment{history}{\comment}{\endcomment}

% For commutative diagrams we use Xy-pic
\usepackage[all]{xy}

% We use 2cell for 2-commutative diagrams.
\xyoption{2cell}
\UseAllTwocells

% We use multicol for the list of chapters between chapters
\usepackage{multicol}

% This is generally recommended for better output
\usepackage{lmodern}
\usepackage[T1]{fontenc}

% For cross-file-references

% Package for hypertext links:
\usepackage{hyperref}

% For any local file, say "hello.tex" you want to link to please

% Theorem environments.
%
\theoremstyle{plain}
\newtheorem{theorem}[subsection]{Theorem}
\newtheorem{proposition}[subsection]{Proposition}
\newtheorem{lemma}[subsection]{Lemma}

\theoremstyle{definition}
\newtheorem{definition}[subsection]{Definition}
\newtheorem{example}[subsection]{Example}
\newtheorem{exercise}[subsection]{Exercise}
\newtheorem{situation}[subsection]{Situation}

\theoremstyle{remark}
\newtheorem{remark}[subsection]{Remark}
\newtheorem{remarks}[subsection]{Remarks}

\numberwithin{equation}{subsection}

% Macros
%
\def\lim{\mathop{\mathrm{lim}}\nolimits}
\def\colim{\mathop{\mathrm{colim}}\nolimits}
\def\Spec{\mathop{\mathrm{Spec}}}
\def\Hom{\mathop{\mathrm{Hom}}\nolimits}
\def\Ext{\mathop{\mathrm{Ext}}\nolimits}
\def\SheafHom{\mathop{\mathcal{H}\!\mathit{om}}\nolimits}
\def\SheafExt{\mathop{\mathcal{E}\!\mathit{xt}}\nolimits}
\def\Sch{\mathit{Sch}}
\def\Mor{\mathop{\mathrm{Mor}}\nolimits}
\def\Ob{\mathop{\mathrm{Ob}}\nolimits}
\def\Sh{\mathop{\mathit{Sh}}\nolimits}
\def\NL{\mathop{N\!L}\nolimits}
\def\CH{\mathop{\mathrm{CH}}\nolimits}
\def\proetale{{pro\text{-}\acute{e}tale}}
\def\etale{{\acute{e}tale}}
\def\QCoh{\mathit{QCoh}}
\def\Ker{\mathop{\mathrm{Ker}}}
\def\Im{\mathop{\mathrm{Im}}}
\def\Coker{\mathop{\mathrm{Coker}}}
\def\Coim{\mathop{\mathrm{Coim}}}

% Boxtimes
%
\DeclareMathSymbol{\boxtimes}{\mathbin}{AMSa}{"02}

%
% Macros for moduli stacks/spaces
%
\def\QCohstack{\mathcal{QC}\!\mathit{oh}}
\def\Cohstack{\mathcal{C}\!\mathit{oh}}
\def\Spacesstack{\mathcal{S}\!\mathit{paces}}
\def\Quotfunctor{\mathrm{Quot}}
\def\Hilbfunctor{\mathrm{Hilb}}
\def\Curvesstack{\mathcal{C}\!\mathit{urves}}
\def\Polarizedstack{\mathcal{P}\!\mathit{olarized}}
\def\Complexesstack{\mathcal{C}\!\mathit{omplexes}}
% \Pic is the operator that assigns to X its picard group, usage \Pic(X)
% \Picardstack_{X/B} denotes the Picard stack of X over B
% \Picardfunctor_{X/B} denotes the Picard functor of X over B
\def\Pic{\mathop{\mathrm{Pic}}\nolimits}
\def\Picardstack{\mathcal{P}\!\mathit{ic}}
\def\Picardfunctor{\mathrm{Pic}}
\def\Deformationcategory{\mathcal{D}\!\mathit{ef}}

\setlength{\emergencystretch}{3em}
\begin{document}
\title[Stacks Project, Tag 08P2]{709 --- Underlying simplicial-set homotopy equivalences induce abelian homology isomorphisms}
\maketitle
\noindent Copyright (C) 2005--2025 Johan de Jong. Stacks Project authors. Source: \href{https://stacks.math.columbia.edu/tag/08P2}{Tag 08P2}.

\noindent Editorial difficulty note (not author text). Difficulty H2: The proof cannot assume the homotopy inverse is additive; it passes to free simplicial abelian groups and constructs normalized cycle lifts to prove both surjectivity and injectivity. This bridge from set-level to additive homotopy is a nonformal derived-homological argument..

\noindent Editorial provenance note (not author text). History: excerpt from revision \texttt{a0444\allowbreak 6e57e\allowbreak c1fbc\allowbreak 252a8\allowbreak 71afc\allowbreak ec775\allowbreak 2fb28\allowbreak 07b14}, simplicial.tex, lines 5884--5950. Reference presentation was adapted; original-author context and core support blocks were retained or added. The mathematical wording is unchanged. Licensed under GNU FDL 1.2 or later; no invariant sections or cover texts. See ../licenses/COPYING. Context blocks reproduce author definitions and supporting results; they are not additional counted problems.

\noindent
There is a second chain complex we can associate to
a simplicial object of $\mathcal{A}$. Recall that by
Lemma \href{https://stacks.math.columbia.edu/tag/017V}{lemma splitting abelian category (Tag 017V)}
any simplicial object $U$ of $\mathcal{A}$
is canonically split with
$N(U_m) = \bigcap_{i = 0}^{m - 1} \Ker(d^m_i)$.
We define the {\it normalized chain complex $N(U)$}
to be the chain complex
$$
\ldots \to N(U_2) \to N(U_1) \to N(U_0) \to 0 \to 0 \to \ldots
$$
with boundary map $d_n : N(U_n) \to N(U_{n - 1})$ given
by the restriction of $(-1)^nd^n_n$ to the direct summand
$N(U_n)$ of $U_n$. Note that Lemma \href{https://stacks.math.columbia.edu/tag/017X}{lemma N d in N (Tag 017X)}
implies that $d^n_n(N(U_n)) \subset N(U_{n - 1})$.
It is a complex because
$d^n_n \circ d^{n + 1}_{n + 1} = d^n_n \circ d^{n + 1}_n$
and $d^{n + 1}_n$ is zero on $N(U_{n + 1})$ by definition.
Thus we obtain a second functor
$$
N : \text{Simp}(\mathcal{A}) \longrightarrow \text{Ch}_{\geq 0}(\mathcal{A}).
$$


\begin{lemma}
\label{lemma-homotopy-equivalence}
Let $f : X \to Y$ be a map of simplicial abelian groups. If $f$ is a
homotopy equivalence of simplicial sets, then $f$ induces a
quasi-isomorphism of associated chain complexes.
\end{lemma}

\begin{proof}
In this proof we will write $H_n(Z) = H_n(s(Z)) = H_n(N(Z))$
when $Z$ is a simplicial abelian group, with $s$ and $N$ as in
Section \href{https://stacks.math.columbia.edu/tag/0194}{section complexes (Tag 0194)}.
Let $\mathbf{Z}[X]$ denote the free abelian group on $X$ viewed as
a simplicial set and similarly for $\mathbf{Z}[Y]$. Consider the
commutative diagram
$$
\xymatrix{
\mathbf{Z}[X] \ar[r]_g \ar[d] &
\mathbf{Z}[Y] \ar[d] \\
X \ar[r]^f & Y
}
$$
of simplicial abelian groups. Since taking the free abelian group on a set
is a functor, we see that the horizontal arrow is a homotopy equivalence of
simplicial abelian groups, see Lemma \href{https://stacks.math.columbia.edu/tag/019Y}{lemma functorial homotopy (Tag 019Y)}.
By Lemma \href{https://stacks.math.columbia.edu/tag/019T}{lemma homotopy equivalence s N (Tag 019T)}
we see that $H_n(g) : H_n(\mathbf{Z}[X]) \to H_n(\mathbf{Z}[Y])$ is bijective
for all $n \geq 0$.

\medskip\noindent
Let $\xi \in H_n(Y)$. By definition of $N(Y)$ we can represent
$\xi$ by an element $y \in N(Y_n)$ whose boundary is zero.
This means $y \in Y_n$ with $d^n_0(y) = \ldots = d^n_{n - 1}(y) = 0$
because $y \in N(Y_n)$ and $d^n_n(y) = 0$ because the boundary
of $y$ is zero. Denote $0_n \in Y_n$ the zero element. Then we see that
$$
\tilde y = [y] - [0_n] \in (\mathbf{Z}[Y])_n
$$
is an element with $d^n_0(\tilde y) = \ldots = d^n_{n - 1}(\tilde y) = 0$
and $d^n_n(\tilde y) = 0$. Thus $\tilde y$ is in $N(\mathbf{Z}[Y])_n$
has boundary $0$, i.e., $\tilde y$ determines a
class $\tilde \xi \in H_n(\mathbf{Z}[Y])$ mapping to $\xi$.
Because $H_n(\mathbf{Z}[X]) \to H_n(\mathbf{Z}[Y])$ is bijective
we can lift $\tilde \xi$ to a class in $H_n(\mathbf{Z}[X])$.
Looking at the commutative diagram above
we see that $\xi$ is in the image of $H_n(X) \to H_n(Y)$.

\medskip\noindent
Let $\xi \in H_n(X)$ be an element mapping to zero in $H_n(Y)$.
Exactly as in the previous parapgraph we can represent
$\xi$ by an element $x \in N(X_n)$ whose boundary is zero, i.e.,
$d^n_0(x) = \ldots = d^n_{n - 1}(x) = d^n_n(x) = 0$.
In particular, we see that $[x] - [0_n]$ is an element of
$N(\mathbf{Z}[X])_n$ whose boundary is zero, whence defines
a lift $\tilde \xi \in H_n(\mathbf{Z}[x])$ of $\xi$.
The fact that $\xi$ maps to zero in $H_n(Y)$
means there exists a $y \in N(Y_{n + 1})$ whose boundary is $f_n(x)$.
This means $d^{n + 1}_0(y) = \ldots = d^{n + 1}_n(y) = 0$ and
$d^{n + 1}_{n + 1}(y) = f(x)$.
However, this means exactly that
$z = [y] - [0_{n + 1}]$ is in $N(\mathbf{Z}[y])_{n + 1}$ and
$$
g([x] - [0_n]) = [f(x)] - [0_n] = \text{boundary of }z
$$
This proves that $\tilde \xi$ maps to zero in $H_n(\mathbf{Z}[y])$.
As $H_n(\mathbf{Z}[X]) \to H_n(\mathbf{Z}[Y])$ is bijective
we conclude $\tilde \xi = 0$ and hence $\xi = 0$.
\end{proof}
\end{document}
